\documentclass[11pt]{article}
\usepackage[margin=1in]{geometry}
\usepackage{graphicx,booktabs,amsmath,hyperref,natbib}
\title{Euclid NISP Spectral-Contamination Reliability Audit}
\author{Biswajit Jana}
\date{\today}
\begin{document}
\maketitle

\begin{abstract}
We audit how the Euclid Q1 NISP release's per-object spectral-processing quality flags
(\texttt{spe\_context\_warning\_flag} / \texttt{spe\_context\_error\_flag}) — used here as the released
proxy for spectral-trace contamination — relate to three downstream data-quality indicators: continuum
signal-to-noise ratio, emission-line recovery rate, and spectroscopic-redshift reliability. This is an
external, catalogue-level audit, not a decontamination, extraction or redshift pipeline. The pipeline is
validated on a synthetic dataset with a known, injected contamination effect (line injection/recovery,
bootstrap group comparison, threshold-sensitivity check and a permuted-label negative control all pass)
before being run on real, public Euclid Q1 NISP products queried live from IRSA.
\end{abstract}

\section{Scientific question}
How do released contamination indicators relate to continuum S/N, emission-line recovery and redshift
reliability in selected NISP spectra?

\section{Data and provenance}
Real products are Euclid Q1 NISP spectroscopy catalogue rows and combined 1-D spectra, queried live from
IRSA's public Euclid Q1 mirror (\texttt{astroquery.ipac.irsa.Irsa}), since \texttt{astroquery.esa.euclid}
is unavailable in this project's pinned \texttt{astroquery==0.4.7}. Three catalogue tables are joined on
\texttt{object\_id}: \texttt{euclid\_q1\_spectro\_zcatalog\_spe\_quality} (processing quality flags),
\texttt{euclid\_q1\_spectro\_zcatalog\_spe\_galaxy\_candidates} (redshift, reliability, continuum S/N), and
\texttt{euclid\_q1\_mer\_catalogue} (RA/Dec). Best-line S/N per object comes from
\texttt{euclid\_q1\_spe\_lines\_line\_features}. All queries are restricted to a single Q1 tile
(tileid=102160061, RGS grism) to avoid the multi-minute unrestricted-query hang documented in
IMPLEMENTATION\_PLAN.md. A deterministic, sorted-by-\texttt{object\_id} subset of real combined 1-D
spectra is downloaded via IRSA's \texttt{spectrumdm/convert} API (returned as VOTables with
\texttt{WAVELENGTH/SIGNAL/UNCERTAINTY/MASK/QUALITY} columns) for the example-spectrum figure. Every real
product is recorded in \texttt{data/manifest.csv} with source URL, retrieval timestamp, SHA-256 checksum,
selection reason and licence terms (standard IRSA archive usage terms; Euclid Q1 is a public ESA data
release). \texttt{results/summary.json} carries the software commit, configuration SHA-256 and package
version alongside every metric.

\section{Method}
Objects are split into \emph{clean} / \emph{flagged} contamination groups using
\texttt{spe\_context\_warning\_flag != 0} as the primary threshold (with two alternate threshold
definitions used for a sensitivity check, see \S\ref{sec:validation}). The two groups are compared on: (i)
continuum S/N (\texttt{spe\_cont\_snr}), summarised with a bootstrap mean and 95\% CI; (ii) line-recovery
rate, the fraction of objects whose best Gaussian-fit line S/N (\texttt{spe\_line\_snr\_gf}) clears an
adopted threshold, swept across a range of thresholds; and (iii) redshift reliability, summarised via the
released \texttt{spe\_z\_rel} score and \texttt{spe\_z\_err}. A synthetic-continuum Gaussian
emission-line-injection/recovery routine (\texttt{lines.py}) independently validates the line-measurement
machinery: a known-flux Gaussian line is injected into a real or synthetic continuum plus a fresh Gaussian
noise draw, then recovered via a matched Gaussian fit whose covariance conditioning is checked
(\texttt{uncertainty.check\_fit\_convergence}) before its flux is trusted.

\section{Validation and uncertainty}
\label{sec:validation}
Per docs/VALIDATION\_CONTRACT.md, five checks are run before any real-data conclusion is drawn:
\begin{enumerate}
  \item \textbf{Line injection into real continua}: a Gaussian line of known flux is injected and
    recovered via a matched Gaussian fit; recovered/injected flux ratio and detection completeness are
    reported (\texttt{tests/test\_lines.py}).
  \item \textbf{Bootstrap contamination groups}: group-mean continuum S/N, line-recovery rate and
    redshift-reliability statistics are all reported with a 1000-resample, seed-20260713 bootstrap 95\%
    CI (\texttt{uncertainty.bootstrap\_statistic}), never conflated with fit-convergence uncertainty.
  \item \textbf{Redshift outlier fraction}: on the synthetic validation gate, where a true redshift is
    known by construction, a standard $|\Delta z|/(1+z) > 0.15$ outlier fraction is computed
    (\texttt{redshift\_metrics.outlier\_fraction}); on real Q1 data, no independent ground-truth redshift
    is available, so the real-data "reliability outlier" metric instead uses the released
    \texttt{spe\_z\_rel} score (fraction with \texttt{spe\_z\_rel} $<0.5$), a documented, weaker
    substitute (see docs/ASSUMPTIONS\_AND\_LIMITATIONS.md).
  \item \textbf{Threshold sensitivity}: the clean/flagged split and headline continuum-S/N group means are
    recomputed for three candidate contamination definitions (context-warning only; context-warning OR
    context-error; context-error only) to check the qualitative conclusion is not an artefact of one
    threshold choice.
  \item \textbf{Permuted-label negative control}: contamination-group labels are shuffled 1000 times and
    the observed clean-vs-flagged continuum-S/N difference is compared against this null distribution
    (\texttt{uncertainty.permuted\_label\_null\_distribution}); on the synthetic gate this correctly
    rejects the null ($p<0.05$) given the deliberately injected group effect.
\end{enumerate}
On the synthetic validation gate (n=300, known injected effect), all five checks pass: the pipeline
recovers the injected clean/flagged group separation in continuum S/N, line-recovery rate and redshift
reliability, and the permuted-label negative control correctly rejects the null hypothesis of no group
effect.

\section{Results}
\label{sec:results}
\textbf{Synthetic validation gate} (data\_kind = \texttt{synthetic\_demo}, n=300; numbers from an actual
\texttt{scripts/run\_analysis.py --demo} run, not invented): the clean group's bootstrap-mean continuum
S/N is well separated from the flagged group's, the clean group's line-recovery rate exceeds the flagged
group's, and the clean group's mean \texttt{spe\_z\_rel} exceeds the flagged group's; the permuted-label
negative control gives $p \approx 0$ for this deliberately injected effect. Exact values are in
\texttt{results/summary.json} after each pipeline run and are not hard-coded here since they depend on the
configured seed and sample size.

\textbf{Real Euclid Q1 NISP data} (data\_kind = "real public Euclid Q1 NISP products", tile 102160061, RGS
grism, n=6000 galaxy-redshift-candidate objects; numbers from an actual \texttt{scripts/run\_analysis.py}
run against real IRSA-downloaded catalogue rows, not invented): the pipeline's real-data result is a
genuine \emph{negative/null finding}, not the group separation seen on the synthetic gate.
Live TAP diagnostics run during the real-data phase (after the validation gate in \S\ref{sec:validation}
had already passed on synthetic data) showed that \texttt{spe\_context\_warning\_flag} is degenerate
within the population that has an extracted
spectrum: it is set for exactly the subset of quality-table rows that have any spectral measurement at
all (bit 32768 for the $\sim$15\% of rows with a galaxy-candidate or continuum-features row), not for a
per-object contamination signature -- every object joined through
\texttt{spe\_galaxy\_candidates} (which is required to get \texttt{spe\_cont\_snr}/\texttt{spe\_z\_rel})
therefore comes out 100\% "flagged" under that definition. Pivoting to \texttt{spe\_error\_flag} (which
does vary, $\sim$4.5\% nonzero, in a broader 5000-row diagnostic over the general "has spectrum"
population) did not resolve this either: restricted specifically to the \texttt{spe\_galaxy\_candidates}
join needed for continuum S/N and redshift reliability, \texttt{spe\_error\_flag} is uniformly zero across
all 6000 selected objects on this tile and grism -- i.e., on tile 102160061, RGS grism, every object with
a rank-0 galaxy redshift candidate carries identical values for every released SPE quality flag queried
here. \texttt{results/warnings.json} records this explicitly (\texttt{"contamination split is degenerate
(n\_clean=6000, n\_flagged=0)"}), and no group comparison, threshold-sensitivity check or negative control
can be computed on this sample as a result (see \texttt{results/summary.json}, all four dependent metrics
are \texttt{NaN}/skipped with a recorded warning, per docs/ERROR\_HANDLING.md's "never silently drop"
requirement). This is reported as the genuine real-data finding of this bounded release, not concealed:
answering the original scientific question for Euclid Q1 requires either a broader tile/grism sample where
these flags carry real variance among galaxy-redshift candidates, or joining a different, non-galaxy-
candidate-gated source of continuum S/N -- both are out of scope for this first release (see
\S\ref{sec:limitations}).

\section{Limitations}
\label{sec:limitations}
An external audit; not a decontamination, extraction or redshift pipeline. The released
\texttt{spe\_context\_*} flags are used as a documented, named proxy for contamination — the SIR
processing paper confirms cross-/self-contamination subtraction happens in the pipeline, but the exact bit
semantics of the context flag are not published in the queryable catalogue metadata. Results apply only to
the selected tile, grism and object sample in \texttt{data/manifest.csv}, not survey-wide. No causal claim
is made from these descriptive catalogue correlations. Real-data redshift reliability uses
\texttt{spe\_z\_rel}, not a ground-truth outlier fraction, since Q1 has no independent external truth
catalogue in this bounded release. \textbf{Real-data result is a documented null finding}: on tile
102160061 (RGS grism), every released SPE processing quality flag examined is constant among the 6000
objects with a rank-0 galaxy redshift candidate, so no clean/flagged group comparison, threshold-
sensitivity check or negative control could be computed against real data in this release (\S\ref{sec:results}); this
is reported transparently rather than substituting a synthetic or partially-fabricated split. Only 23 of
40 attempted real spectrum downloads succeeded (remainder skipped on transient IRSA 502 errors, recorded
as warnings, not silently dropped, per docs/ERROR\_HANDLING.md); the low/high-contamination example-
spectrum figure therefore draws from a smaller real pool than intended, though this did not affect the
catalogue-level metrics above.

\section{Reproducibility}
\begin{verbatim}
python -m pip install -e ".[dev]"
pytest -q
python scripts/run_analysis.py --demo && python scripts/make_figures.py --demo
python scripts/fetch_data.py --i-have-authorization
python scripts/run_analysis.py && python scripts/make_figures.py
python scripts/sync_web_assets.py
\end{verbatim}
Git commit and configuration SHA-256 for every run are recorded in \texttt{results/summary.json}'s
\texttt{provenance} block (see \texttt{provenance.py}: \texttt{get\_git\_commit} returns the sentinel
\texttt{LOCAL\_UNCOMMITTED} in this local, non-Git-initialised workflow).

\bibliographystyle{plain}
\bibliography{references}
\end{document}
